%===========================================================================


\section{External Validation on eICU-CRD \texorpdfstring{[C]}{[C]}}
\label{sec:external_validation}
\label{sec:external}
%===========================================================================

External validation asks whether a frozen model ranks person-hours as well in a
different health system as in the one it was fitted on. Because the Sepsis-3
label must be transported as well as the model, this section reports the
labelling inputs first, then discrimination, then operating characteristics.

\subsection{Coverage of the Sepsis-3 Anchor in eICU-CRD}
\label{sec:external_coverage}

In eICU-CRD the culture limb of the Sepsis-3 anchor is documented for
\pcExtMicroCoverage\% of the analysis cohort (\pcExtNMicroStays{} of
\pcExtNCohortStays{} unit stays), whereas antibiotic exposure, pooled across
the \texttt{medication} and \texttt{infusionDrug} tables, reaches
\pcExtAbxCoverage\% of stays. Both limbs are documented at any offset for
\pcExtNBothLimbStays{} stays (\pcExtBothLimbCoverage\%); they co-occur inside
each variant's antibiotic--culture window for \pcExtNSuspicionA{},
\pcExtNSuspicionB{} and \pcExtNSuspicionC{} stays under Variants A, B and C;
and adding the organ-dysfunction criterion leaves \pcExtNOnsetStaysA{},
\pcExtNOnsetStaysB{} and \pcExtNOnsetStaysC{} labelled stays.

Under Variant B, \pcExtNOnsetStaysB{} stays are labelled but \pcExtNEventsB{}
events are scored: an onset after a stay's last observed hour is treated as
right-censoring, the horizon rule the development cohort applies at 72 hours.
Two of Variant B's onsets fall this way and none of Variant A's or C's; every
other labelled onset hour receives a prediction from all three models.

Because no onset is observable in a stay without a culture record, the primary
external analysis is restricted to the \pcExtPanelHoursMicro{} person-hours of
the microbiology-covered sub-cohort rather than all \pcExtPanelHoursFull{}.
Table~\ref{tab:external_rates} reconciles the resulting event rates with the
development cohort's.

\begin{table}[H]
    \centering\small
    \adjustbox{max width=\textwidth}{%
    \begin{tabular}{llrrr}
        \toprule
        \textbf{Var} & \textbf{Cohort} & \textbf{Person-hours}
        & \textbf{Events} & \textbf{Events / 1{,}000 person-h} \\
        \midrule
        A & MIMIC-IV temporal test        & \pcNTestRowsA & \pcNTestEventsA & \pcExtRateMimicA \\
        A & eICU, microbiology-covered    & \pcExtNRowsA  & \pcExtNEventsA  & \pcExtRateEicuA \\
        A & eICU, antibiotic-only anchor  & \pcExtAbxNRowsA & \pcExtAbxNEventsA & \pcExtRateEicuAbxA \\
        \midrule
        B & MIMIC-IV temporal test        & \pcNTestRowsB & \pcNTestEventsB & \pcExtRateMimicB \\
        B & eICU, microbiology-covered    & \pcExtNRowsB  & \pcExtNEventsB  & \pcExtRateEicuB \\
        B & eICU, antibiotic-only anchor  & \pcExtAbxNRowsB & \pcExtAbxNEventsB & \pcExtRateEicuAbxB \\
        \midrule
        C & MIMIC-IV temporal test        & \pcNTestRowsC & \pcNTestEventsC & \pcExtRateMimicC \\
        C & eICU, microbiology-covered    & \pcExtNRowsC  & \pcExtNEventsC  & \pcExtRateEicuC \\
        C & eICU, antibiotic-only anchor  & \pcExtAbxNRowsC & \pcExtAbxNEventsC & \pcExtRateEicuAbxC \\
        \bottomrule
    \end{tabular}}
    \caption{Event-rate reconciliation for the development cohort and both
    external anchors. The antibiotic-only anchor is not Sepsis-3 and its rate
    is not interchangeable with the other two.}
    \label{tab:external_rates}
\end{table}

On the restricted cohort the external event rate is \pcExtRateEicuB{} per
thousand person-hours against MIMIC-IV's \pcExtRateMimicB{} under the primary
label, a ratio of \pcExtRateRatioB. Relaxing the culture requirement raises it
to \pcExtRateEicuAbxB{} per thousand (ratio \pcExtRateRatioAbxB), suggesting
that much of the gap lies in the culture limb. The residual includes case mix,
charting frequency and differences in illness severity;
Section~\ref{sec:external_altlabel} narrows it further by removing the anchor
altogether. A gap of two or more orders of magnitude between two large adult
ICU populations would be far more consistent with missing labelling inputs
than with a difference in incidence, and would call for an audit of those
inputs before the rate is interpreted.

\subsection{Discrimination Under Transport}

\begin{table}[H]
    \centering
    \adjustbox{max width=\textwidth}{%
    \begin{tabular}{llrrrr}
        \toprule
        \textbf{Var} & \textbf{Model} & \textbf{N rows} & \textbf{N events}
        & \textbf{AUROC (eICU)} & \textbf{AUROC (MIMIC)} \\
        \midrule
        A & Primary & \pcExtNRowsA & \pcExtNEventsA & \pcExtAurocPrimaryA & \pcAurocPrimaryA \\
        A & GBT     & \pcExtNRowsA & \pcExtNEventsA & \pcExtAurocGbtA  & \pcAurocGbtA \\
        A & NEWS2   & \pcExtNRowsA & \pcExtNEventsA & \pcExtAurocNewsA & \pcAurocNewsA \\
        \midrule
        B & Primary & \pcExtNRowsB & \pcExtNEventsB & \pcExtAurocPrimaryB & \pcAurocPrimaryB \\
        B & GBT     & \pcExtNRowsB & \pcExtNEventsB & \pcExtAurocGbtB  & \pcAurocGbtB \\
        B & NEWS2   & \pcExtNRowsB & \pcExtNEventsB & \pcExtAurocNewsB & \pcAurocNewsB \\
        \midrule
        C & Primary & \pcExtNRowsC & \pcExtNEventsC & \pcExtAurocPrimaryC & \pcAurocPrimaryC \\
        C & GBT     & \pcExtNRowsC & \pcExtNEventsC & \pcExtAurocGbtC  & \pcAurocGbtC \\
        C & NEWS2   & \pcExtNRowsC & \pcExtNEventsC & \pcExtAurocNewsC & \pcAurocNewsC \\
        \bottomrule
    \end{tabular}}
    \caption{External validation on eICU-CRD (frozen models, no retraining),
    Sepsis-3 anchor on the microbiology-covered sub-cohort. Each variant
    carries its own external label set, so the row counts and event counts
    differ across variants as they do internally.}
    \label{tab:external}
\end{table}

\textbf{The point estimate falls under every label.} The primary
model falls from \pcAurocPrimaryA{} to \pcExtAurocPrimaryA{} (Narrow), from
\pcAurocPrimaryB{} to \pcExtAurocPrimaryB{} (Seymour-Standard) and from
\pcAurocPrimaryC{} to \pcExtAurocPrimaryC{} (Liberal). Only the
Seymour-Standard contrast, the tested contrast of H5, carries an interval, and
it covers zero (Section~\ref{sec:h5}); the Narrow and Liberal figures are point
estimates on \pcExtNEventsA{} and \pcExtNEventsC{} events, so the three
indicate a direction rather than three established degradations. The
Seymour-Standard drop of \pcHFiveEst{} is about three-quarters of the 0.085
reference estimate reported by Moor et al.\ \cite{moor2023review}. The point
estimates show the same cross-label ordering as the internal results, although
the low event counts preclude interpreting that ordering.

NEWS2 falls to \pcExtAurocNewsB{} under the primary label, close to chance. On
this arm the fitted models therefore retain more of their advantage over a
bedside score, although under the antibiotic-only anchor, with two orders of
magnitude more events, the ordering between the boosted model and NEWS2 does
not hold (Section~\ref{sec:external_interpretability}).

The differing event counts and NEWS2 AUROCs across variants are consistent
with variant-specific external labels: NEWS2's score is label-invariant, but
its AUROC depends on the outcome it is ranked against.

Every external event count above falls below the interpretability floor
(Section~\ref{sec:external_interpretability}).

\subsection{The Anchor-Independent Label on the Full External Cohort
\texorpdfstring{[D]}{[D]}}
\label{sec:external_altlabel}

The estimates above are computed on the \pcExtNMicroStays{} stays
(\pcExtMicroCoverage\%) with documented microbiology. Variant D carries no
treatment timestamp, so it can be constructed on the whole cohort. This
subsection reports the frozen Variant D models applied to all
\pcExtAltNStays{} eICU-CRD stays under Protocol Amendment~9.

\begin{table}[H]
    \centering\small
    \begin{tabular}{lrrlr}
        \toprule
        \textbf{Model} & \textbf{Internal AUROC} & \textbf{External AUROC}
        & \textbf{Drop (95\% CI)} & \textbf{Events} \\
        \midrule
        Primary & \pcAltAurocFullPrimaryD & \pcExtAltAurocPrimary & \pcExtAltDropPrimary{} \pcExtAltDropCiPrimary & \pcExtAltNEvents \\
        GBT     & \pcAltAurocFullGbtD     & \pcExtAltAurocGbt     & \pcExtAltDropGbt{} \pcExtAltDropCiGbt & \pcExtAltNEvents \\
        NEWS2   & \pcNotApplicable        & \pcExtAltAurocNews    & \pcNotApplicable & \pcExtAltNEvents \\
        \bottomrule
    \end{tabular}
    \caption[Variant D on the full eICU-CRD cohort]{Variant D applied to all
    \pcExtAltNStays{} eICU-CRD stays (\pcExtAltNRows{} person-hours). The target
    is acute organ dysfunction, not Sepsis-3, so these values are not pooled
    with the Sepsis-3 arm. eICU-CRD's periodic vitals carry no Glasgow Coma
    Scale, so the external SOFA delta rests on five components against
    MIMIC-IV's six. The drop's interval resamples the two databases
    independently and differences the streams, as for H5 (Protocol
    Amendment~10). Post-hoc under Protocol Amendment~9; excluded from H5.}
    \label{tab:external_altlabel}
\end{table}

The arm scores \pcExtAltNEvents{} events over \pcExtAltNRows{} person-hours,
against \pcExtNEventsB{} under Variant B, and is the only external estimate that
clears the \pcExtMinEvents-event interpretability floor
(Section~\ref{sec:external_interpretability}).

\textbf{AUROC drop.} The primary
model falls from \pcAltAurocFullPrimaryD{} internally to
\pcExtAltAurocPrimary{} externally, a difference of
\pcExtAltDropPrimary{} (95\% CI \pcExtAltDropCiPrimary), and the
gradient-boosted comparator from
\pcAltAurocFullGbtD{} to \pcExtAltAurocGbt, a difference of
\pcExtAltDropGbt{} (\pcExtAltDropCiGbt). Both drops are of the order of the
0.085 degradation Moor et al.\ \cite{moor2023review} report across four
databases, and of the same order as the H5 estimate. The event count affects
the width rather than the location: both intervals exclude zero and are
narrower than a hundredth of a point, whereas H5's substituted contrast,
\pcHFiveEst{} [\pcHFiveCiLo, \pcHFiveCiHi] on \pcExtNEventsB{} events, is
compatible with both no degradation and a degradation twice as large. For the
deterioration target, transport between these databases therefore reduces
discrimination; this is not evidence about the Sepsis-3 target, a different
estimand.

\textbf{Event rate.} Under the Sepsis-3 anchor the external event rate is
\pcExtRateRatioB{} of the development cohort's, a gap of roughly an order of
magnitude located largely in the culture limb
(Section~\ref{sec:external_coverage}). Under the anchor-independent label the
two databases give \pcExtAltRateMimic{} and \pcExtAltRateEicu{} events per
1{,}000 person-hours, a ratio of \pcExtAltRateRatio. The reduction in the
event-rate gap after removing the treatment anchor suggests that label
transport accounts for a substantial part of the original discrepancy
(Section~\ref{sec:variant_d_establishes}).

\subsection{External Clinical Utility on eICU-CRD \texorpdfstring{[C]}{[C]}}

Operating characteristics are reported on the same microbiology-covered
sub-cohort and inherit its coverage limitation
(Section~\ref{sec:external_coverage}).

\begin{table}[H]
    \centering\small
    \begin{tabular}{llrrrrr}
        \toprule
        \textbf{Var} & \textbf{Model} & \textbf{Specificity} & \textbf{PPV}
        & \textbf{NNE} & \textbf{FA/pt-day} & \textbf{Patient PPV} \\
        \midrule
        A & Primary & \pcExtSpecPrimaryA & \pcExtPpvPrimaryA & \pcExtNnePrimaryA & \pcExtFaPerPtDayPrimaryA & \pcExtPatientPpvPrimaryA \\
        B & Primary & \pcExtSpecPrimaryB & \pcExtPpvPrimaryB & \pcExtNnePrimaryB & \pcExtFaPerPtDayPrimaryB & \pcExtPatientPpvPrimaryB \\
        C & Primary & \pcExtSpecPrimaryC & \pcExtPpvPrimaryC & \pcExtNnePrimaryC & \pcExtFaPerPtDayPrimaryC & \pcExtPatientPpvPrimaryC \\
        \midrule
        A--C & NEWS2 & \pcExtSpecNewsB & \pcExtPpvNewsB & \pcExtNneNewsB & \pcExtFaPerPtDayNewsB & \pcExtPatientPpvNewsB \\
        \bottomrule
    \end{tabular}
    \caption{eICU-CRD operating characteristics at the same 80\% target
    sensitivity, thresholds re-derived on the external cohort. All rows are
    drawn from the same \pcExtNStays{} microbiology-covered stays, but each
    variant carries its own label and therefore its own denominator:
    \pcExtNRowsA, \pcExtNRowsB{} and \pcExtNRowsC{} person-hours on
    \pcExtNEventsA, \pcExtNEventsB{} and \pcExtNEventsC{} event hours under
    A, B and C (Table~\ref{tab:external}). The NEWS2 row is scored under
    Variant B. Re-deriving thresholds externally excludes the further loss a
    threshold transported from MIMIC-IV would incur.}
    \label{tab:external_utility}
\end{table}

Operating characteristics on eICU-CRD are substantially worse than internal ones
at the same target sensitivity. Reaching 80 per cent sensitivity requires
\pcExtNnePrimaryC{} to \pcExtNnePrimaryA{} alerted person-hours per true event
hour, at \pcExtFaPerPtDayPrimaryB--\pcExtFaPerPtDayPrimaryA{} false-alert hours
per patient-day, roughly one alert every two hours per patient. NEWS2 cannot
reach the target without alerting on nearly
every person-hour (specificity \pcExtSpecNewsB, \pcExtFaPerPtDayNewsB{}
false-alert hours per patient-day). Patient-level PPV of
\pcExtPatientPpvPrimaryB{} is again the cohort septic fraction
(\pcExtNSepsisStays/\pcExtNStays), because the model alerts on essentially
every stay. At this event count the contributions of model and label cannot be
separated, since PPV, NNE and false-alarm rate all scale with the event rate;
these figures are therefore reported as an upper bound on transferability and
support no separate conclusion.
